% ECONOMICS_PROBLEM: {"id": "H21-85", "title": "Optimal income taxation with heterogeneous people", "jel_code": "H21", "source_id": "OP-0470"}
% !TeX program = xelatex
\documentclass[11pt,a4paper]{article}
\usepackage[margin=23mm]{geometry}
\usepackage{fontspec}
\setmainfont{Latin Modern Roman}
\usepackage{amsmath,amssymb,mathtools}
\usepackage{xcolor,enumitem,booktabs,longtable,array,xurl,hyperref}
\definecolor{muted}{HTML}{53636C}
\hypersetup{colorlinks=true,urlcolor=blue,linkcolor=blue}
\setlength{\parindent}{0pt}
\setlength{\parskip}{5pt}
\newcommand{\R}{\mathbb R}
\newcommand{\E}{\mathbb E}
\newcommand{\N}{\mathbb N}
\newcommand{\ind}{\mathbf 1}
\newcommand{\Var}{\operatorname{Var}}
\newcommand{\Cov}{\operatorname{Cov}}
\newcommand{\supp}{\operatorname{supp}}
\DeclareMathOperator*{\argmin}{arg\,min}
\DeclareMathOperator*{\argmax}{arg\,max}
\newcommand{\entrymeta}[3]{{\sffamily\small\color{muted}\textbf{\texttt{#1}}\quad|\quad #2\par #3\par}}
\newcommand{\Problemref}[1]{\texttt{#1}}
\newcommand{\JELref}[2]{\texttt{#2} (JEL \texttt{#1})}
\begin{document}
\section*{Optimal income taxation with heterogeneous people}
\textbf{Problem ID:} \texttt{H21-85}\quad \textbf{JEL:} \texttt{H21}\par
% BEGIN ECONOMICS STATEMENT
\label{problem:OP-0470}
{\sffamily\small\color{muted}Original subject group: Taxation and social insurance\par}
\label{definitions:problem:30007752}
\textbf{Audit classification:} Published research agenda. The income-tax agenda is verified; the minimax formulation is editorial. The original omitted a complete resource constraint, global household tie handling, and cardinal welfare normalization.
\subsection*{Definitions and precise research target}
Fix a Borel probability space of types \((\Theta,F)\), finitely many occupations \(j=1,\ldots,J\), labor bound \(\bar\ell<\infty\), bounded measurable abilities \(a_j(\theta)>0\), and continuous utility indices \(u_\theta(c,\ell,j)\), increasing in consumption and decreasing in labor. The cardinal normalization of each utility and the social functions \(G_\theta\) are primitives: interpersonal welfare comparisons cannot be inferred from ordinal choices alone. Let \(f:\mathbb R_+^J\to\mathbb R_+\) be continuous and concave on the closed cone, homogeneous of degree one, and continuously differentiable on its positive interior. At a boundary labor vector use finite supporting wages, when they exist; differentiability at the origin is not imposed. An admissible tax schedule \(T:\mathbb R_+\to\mathbb R\) is Lipschitz with declared slope and intercept bounds. Negative taxes are transfers. Fix a government purchase requirement \(R\geq0\), a disposal rule for any surplus receipts, and a price normalization in units of output.

An equilibrium \(e\in\mathcal E(T;F)\) consists of wages \(w\), a measurable probability kernel \(\kappa_\theta^e(d\ell,dj)\) of household choices, consumption and effective labor totals \(L_j\). Every choice in a type's support globally maximizes \(u_\theta(w_ja_j(\theta)\ell-T(w_ja_j(\theta)\ell),\ell,j)\) over \(0\leq\ell\leq\bar\ell\), occupations and nonnegative consumption. Effective labor is the integral of \(a_j(\theta)\ell\) supplied in occupation \(j\); firms maximize profits and \(w_j=\partial_j f(L)\) when differentiable and interior. Household consumption plus public purchases and declared disposal equals output. Restrict to taxes with nonempty equilibria and bounded integrable welfare. Define \(u_\theta(e)=\int u_\theta(c,\ell,j)\,d\kappa_\theta^e\) and \(r_\theta(e)=\int T(w_ja_j(\theta)\ell)\,d\kappa_\theta^e\), with consumption given by the tax budget. Resource and labor aggregates likewise integrate over this kernel. Define the robustly feasible tax class
\[
 \mathcal T_F=\left\{T\in\mathcal T:\ \mathcal E(T;F)\ne\varnothing,\quad
 \int r_\theta(e)\,dF(\theta)\ge R\ \text{for all }e\in\mathcal E(T;F)\right\}.
\]
The robust welfare value is
\[
 V(F)=\sup_{T\in\mathcal T_F}\inf_{e\in\mathcal E(T;F)}
 \int G_\theta(u_\theta(e))\,dF(\theta).
\]
An empty feasible set must be reported; an optimal schedule is claimed only if the supremum is attained. Characterize welfare-improving reforms or \(\varepsilon\)-optimal schedules for specified primitives, and determine which features of the joint ability/preference distribution are identified by the declared observation model. A household bargaining weight alone is insufficient: household utilities, feasible intrafamily allocations and the allocation rule must be specified before treating the household as a type. This is an editorial extension of a published agenda, not a universal optimal-tax theorem.
\subsection*{Variants and assumption changes}
\begin{enumerate}
\item \textbf{Fixed wages versus equilibrium wages (editorial).} Replace production by exogenous \(w\). Compare this partial-equilibrium value with \(V(F)\); a reform can alter occupations and relative wages in the second model.
\item \textbf{Identified preferences versus ambiguity (editorial).} Given a data law \(P_O\), define \(\mathcal F(P_O)\) as distributions generating that law under the fixed observation model. Optimize against all \(F\in\mathcal F(P_O)\), keeping the cardinal welfare normalization fixed; ordinal observational equivalence does not identify welfare weights.
\item \textbf{Individuals versus households (editorial).} Replace \(u_\theta\) by a specified collective-household model with individual consumption, labor and policy-dependent bargaining. This changes both behavior and whose utility enters welfare.
\end{enumerate}
\subsection*{References and source audit}
The 2022 Harvard manuscript, sections 3.2, 8.2 and 9 (printed pp. 89--90), supports the heterogeneity/endogenous-wage agenda. The AEA record confirms the 2024 journal version and DOI. Neither source states this minimax model.
\begin{enumerate}\raggedright
\item\hypertarget{definitions:source:30007752:1}{} Louis Kaplow. 2022. \emph{Optimal Income Taxation}. Olin Discussion Paper 1083, July 2022; sections 3.2, 7, 8.2 and 9, printed pp. 89–90.
\url{https://laweconcenter.law.harvard.edu/wp-content/uploads/2024/11/Kaplow_1083.pdf}.
DOI: \url{https://doi.org/10.3386/w30199}.
\item\hypertarget{definitions:source:30007752:2}{} Louis Kaplow. 2024. \emph{Optimal Income Taxation}. Journal of Economic Literature 62(2), 637–738; publisher abstract.
\url{https://www.aeaweb.org/articles?id=10.1257/jel.20221647}.
DOI: \url{https://doi.org/10.1257/jel.20221647}.
\end{enumerate}

\subsection*{Common conventions: Source mathematical conventions}

These conventions apply to each entry unless explicitly replaced.

\textbf{Models and identification.} A primitive model \(m\) specifies all probability laws, feasible choices, information, equilibrium and measurement rules used by its target. The admissible class \(\mathcal M\) is fixed independently of outcomes. If \(P_O(m)\) is its observable probability law and \(\tau(m)\) the target, its identified set at observable law \(P\) is
\[
 \mathcal I_\tau(P)=\{\tau(m):m\in\mathcal M,\ P_O(m)=P\}.
\]
Point identification means that this set is a singleton. An empty set signals incompatibility of data and assumptions. Coordinate bounds need not describe the joint set. Bounds are sharp only if every retained target is attainable or arbitrarily approachable by compatible models. Finite bounds require explicit restrictions on unobserved quantities; unrestricted confounding, hidden wealth, extrapolation or arbitrary equilibrium selection can yield uninformative sets.

\textbf{Causal interventions.} A potential outcome \(Y_i(p)\) is the outcome under a complete policy regime \(p\), including timing, financing, information and market spillovers. The target population and units are fixed before treatment. With interference, use the complete assignment vector or a justified exposure mapping; individual no-interference notation alone cannot identify an economy-wide intervention. A difference of mean potential outcomes does not require the cross-world joint distribution. Conditioning on post-treatment migration, survival, enrollment or employment changes the estimand and needs separate justification.

\textbf{Statistical statements.} An estimator, confidence set, forecast or test specifies a sampling law, model class and loss/coverage criterion. Uniform \((1-\alpha)\)-coverage means \(\inf_{m\in\mathcal M}\Pr_m\{\tau(m)\in C_n\}\geq1-\alpha\), or an explicitly asymptotic version. The whole parameter space trivially has coverage, so informativeness requires a stated width, risk or power objective. A forecasting improvement is relative to a named benchmark, information set and evaluation population.

\textbf{Equilibrium and optimization.} An equilibrium comprises feasible plans, optimality, beliefs where needed, budget constraints and all market/resource clearing. The primitives or a stated benchmark define each correspondence \(\mathcal E(p;m)\). Nonemptiness is assumed or proved, never inferred just from compact policy sets. A selected-equilibrium target uses a declared selection rule; a robust target quantifies over all permitted equilibria. Use suprema or infima unless attainment is established. An \(\varepsilon\)-optimal policy is feasible and within \(\varepsilon\) of the finite optimum value. Report an empty feasible set separately.

\textbf{Welfare and accounting.} Normative utility, interpersonal normalization, social weights, numeraire, discounting and the use of public receipts are primitives. Behavioral utility need not coincide with the welfare criterion. Household welfare includes ownership income and rents once; adding profits again to compensating variation generally double-counts them. A monetary equivalent is defined by an explicit transfer and price convention, with monotonicity and range conditions. Average effects, mechanism contributions, transfers and welfare losses are different targets. Infinite sums require convergence and dynamic feasible plans require terminal or no-Ponzi conditions.

\textbf{Source versions.} A working-paper date, revision date and journal publication year are distinguished. A DOI for a journal version is not automatically the DOI of an earlier manuscript. Locators refer to the stated version and printed pages where specified. A metadata-only check does not verify a claim about a full-text conclusion. Every entry's source subsection narrows the provenance claim accordingly.
% END ECONOMICS STATEMENT
\end{document}
